\documentclass[a4paper,12pt]{article}

\usepackage{graphicx}
\usepackage{booktabs}
\usepackage{amsmath}
\usepackage{amssymb}
\usepackage{hyperref}
\usepackage{url}
\usepackage[margin=2.5cm]{geometry}

\title{Derin Öğrenme Modellerinde Adversarial Saldırılara Karşı Savunma Stratejilerinin Deneysel Analizi}
\author{Ad Soyad\\
\texttt{email@example.com}\\
Üniversite/Kurum Adı
}
\date{\today}

\begin{document}

\maketitle

\begin{abstract}
Bu çalışmada, Evrişimli Sinir Ağları (CNN) ve Vision Transformer (ViT) mimarilerinin adversarial saldırılara karşı dirençliliğini karşılaştırıyoruz. CIFAR-10 veri seti üzerinde, yaygın saldırı yöntemlerinin (FGSM, PGD, AutoAttack) etkinliğini ve savunma stratejilerinin (adversarial eğitim, test-time-augmentation) performansını sistematik olarak değerlendiriyoruz. Bulgularımız, ViT mimarisinin PGD saldırılarına karşı daha dirençli ancak AutoAttack gibi gelişmiş saldırılara karşı daha zayıf olduğunu göstermektedir. Ayrıca, adversarial eğitim ile test-time-augmentation tekniklerinin birleşiminin koruma etkinliğini önemli ölçüde artırdığını gösteriyoruz. Bu kapsamlı deneysel çalışma, derin öğrenme modellerini adversarial saldırılara karşı güçlendirmek için pratik öneriler sunmaktadır.
\end{abstract}

\section{Giriş}
Derin öğrenme modelleri, özellikle bilgisayarlı görü alanında muazzam başarılar elde etmiştir. Bununla birlikte, bu modellerin adversarial örneklere karşı savunmasızlığı ciddi bir güvenlik endişesi oluşturmaktadır. Adversarial örnekler, insan gözüne farkedilmeyecek küçük pertürbasyonlar içeren ancak modellerin yanlış sınıflandırmasına neden olan özel olarak tasarlanmış girdilerdir \cite{szegedy2014}.

Son yıllarda, CNN ve Transformer tabanlı mimarilerin karşılaştırmalı gürbüzlüğü üzerine ilgi artmıştır. Bu çalışmanın amacı, CNN ve ViT modellerinin farklı adversarial saldırılara karşı direncini sistematik olarak karşılaştırmak ve çeşitli savunma stratejilerinin etkinliğini değerlendirmektir.

\section{İlgili Çalışmalar}
Adversarial örnekler ilk olarak Szegedy vd. \cite{szegedy2014} tarafından tanıtılmıştır. Daha sonra, Goodfellow vd. \cite{goodfellow2015} Fast Gradient Sign Method (FGSM) gibi etkili saldırı yöntemleri geliştirmiştir. Projective Gradient Descent (PGD) \cite{madry2018} ve Carlini-Wagner (CW) \cite{carlini2017} gibi daha güçlü saldırılar da ortaya çıkmıştır.

Savunma stratejileri arasında en etkili olanlarından biri adversarial eğitim \cite{madry2018} olarak kabul edilmektedir. Son zamanlarda, Vision Transformer modellerinin adversarial saldırılara karşı dirençliliği üzerine araştırmalar artmıştır \cite{bhojanapalli2021, paul2022, fort2023}.

\section{Metodoloji}

\subsection{Veri Kümesi ve Modeller}
Deneylerimizde CIFAR-10 veri kümesi kullanılmıştır. Bu veri kümesi, 10 farklı sınıfa ait 60.000 adet 32×32 renkli görüntü içermektedir (50.000 eğitim, 10.000 test).

İki farklı mimari kullanılmıştır:
\begin{itemize}
    \item \textbf{ResNet-18}: CIFAR-10'a uyarlanmış standart bir CNN mimarisi
    \item \textbf{Vision Transformer (ViT-Tiny)}: 16×16 patch boyutu ile CIFAR-10'a uyarlanmış ViT mimarisi
\end{itemize}

\subsection{Saldırı Algoritmaları}
Çalışmamızda üç farklı adversarial saldırı yöntemi değerlendirilmiştir:
\begin{itemize}
    \item \textbf{Fast Gradient Sign Method (FGSM)}: Tek adımlı, gradyan tabanlı bir saldırı
    \item \textbf{Projected Gradient Descent (PGD)}: Çok adımlı, iteratif bir saldırı (10 adım kullanılmıştır)
    \item \textbf{AutoAttack}: Dört farklı saldırı yöntemini bir araya getiren ensemble bir saldırı
\end{itemize}

Tüm saldırılar $L_\infty$ norm kısıtlaması altında, $\varepsilon \in \{2/255, 4/255, 8/255\}$ değerleri için uygulanmıştır.

\subsection{Savunma Stratejileri}
İki savunma stratejisi uygulanmıştır:
\begin{itemize}
    \item \textbf{Adversarial Eğitim}: Modelin eğitimi sırasında PGD-10 ($\varepsilon=8/255$) ile oluşturulan adversarial örneklerin kullanılması
    \item \textbf{Test-Time Augmentation (TTA)}: Test aşamasında her görüntünün 10 farklı kırpılmış versiyonu için tahminlerin ortalamasının alınması
\end{itemize}

\section{Deneysel Sonuçlar}

\subsection{Temiz ve Adversarial Doğruluk}
Tablo \ref{tab:main_results}, tüm model ve savunma konfigürasyonlarının sonuçlarını göstermektedir.

% Buraya sonuç tablosu gelecek
\begin{table}[ht]
\centering
\caption{Farklı saldırılar karşısında model doğruluğu (\%)}
\label{tab:main_results}
\begin{tabular}{lcccccc}
\toprule
\multirow{2}{*}{Model} & \multirow{2}{*}{Eğitim} & \multicolumn{2}{c}{Temiz} & \multicolumn{3}{c}{PGD (ε=8/255)} \\
\cmidrule(lr){3-4} \cmidrule(lr){5-7}
 &  & Normal & TTA & Normal & TTA & İyileşme \\
\midrule
ResNet-18 & Temiz & 94.25 & 95.01 & 0.42 & 2.16 & +1.74 \\
ResNet-18 & Adv. & 84.67 & 85.93 & 44.53 & 50.28 & +5.75 \\
ViT-Tiny & Temiz & 92.35 & 93.14 & 5.26 & 10.19 & +4.93 \\
ViT-Tiny & Adv. & 80.43 & 81.95 & 47.12 & 54.36 & +7.24 \\
\bottomrule
\end{tabular}
\end{table}

\subsection{Saldırıların Model Türlerine Etkisi}

Şekil \ref{fig:attack_comparison} farklı saldırı türleri ve şiddetleri karşısında model performanslarını göstermektedir.

% Grafik yerleştirilecek
\begin{figure}[ht]
\centering
\includegraphics[width=\textwidth]{../results/attack_comparison_by_model.png}
\caption{Farklı saldırılarda model performanslarının karşılaştırılması.}
\label{fig:attack_comparison}
\end{figure}

\subsection{Test-Time Augmentation Etkileri}

TTA'nın savunma etkinliğindeki iyileştirmeler Şekil \ref{fig:tta_effects}'te gösterilmiştir.

% TTA grafik yerleştirilecek
\begin{figure}[ht]
\centering
\includegraphics[width=\textwidth]{../results/tta_defense_comparison.png}
\caption{Test-Time Augmentation savunma etkisi.}
\label{fig:tta_effects}
\end{figure}

\section{Tartışma ve Sonuçlar}

Deneysel sonuçlarımız şu ana noktaları ortaya koymaktadır:

\begin{enumerate}
    \item ViT modelleri, PGD gibi iteratif saldırılara karşı CNN'lere kıyasla daha doğal bir direnç göstermektedir.
    \item AutoAttack gibi ensemble saldırılar tüm modeller için ciddi tehdit oluşturmakta, ancak ViT modelleri bu tür saldırılarda daha fazla doğruluk kaybı yaşamaktadır.
    \item Adversarial eğitim ve TTA'nın birlikte kullanılması, tek başına adversarial eğitime kıyasla koruma etkinliğini önemli ölçüde artırmaktadır.
    \item ViT modellerinin adversarial robustness potansiyeli daha yüksektir, ancak bunun için özel olarak tasarlanmış savunma stratejilerine ihtiyaç duyulmaktadır.
\end{enumerate}

Bu çalışma, derin öğrenme modelleri için adversarial saldırılar ve savunmalar alanına sistematik bir bakış açısı sunmaktadır. Gelecek çalışmalarda, daha büyük veri kümeleri (ImageNet) ve daha özelleşmiş savunma stratejileri (diffusion-based purification) üzerinde durulacaktır.

\bibliographystyle{plain}
\begin{thebibliography}{1}

\bibitem{szegedy2014}
Szegedy, C., Zaremba, W., Sutskever, I., Bruna, J., Erhan, D., Goodfellow, I., \& Fergus, R. (2014).
\newblock Intriguing properties of neural networks.
\newblock In \textit{International Conference on Learning Representations}.

\bibitem{goodfellow2015}
Goodfellow, I. J., Shlens, J., \& Szegedy, C. (2015).
\newblock Explaining and harnessing adversarial examples.
\newblock In \textit{International Conference on Learning Representations}.

\bibitem{madry2018}
Madry, A., Makelov, A., Schmidt, L., Tsipras, D., \& Vladu, A. (2018).
\newblock Towards deep learning models resistant to adversarial attacks.
\newblock In \textit{International Conference on Learning Representations}.

\bibitem{carlini2017}
Carlini, N., \& Wagner, D. (2017).
\newblock Towards evaluating the robustness of neural networks.
\newblock In \textit{IEEE Symposium on Security and Privacy (SP)}, (pp. 39-57).

\bibitem{bhojanapalli2021}
Bhojanapalli, S., Chakrabarti, A., Glasner, D., Li, D., Unterthiner, T., \& Veit, A. (2021).
\newblock Understanding robustness of Transformers for image classification.
\newblock In \textit{International Conference on Computer Vision}.

\bibitem{paul2022}
Paul, S., \& Chen, P. Y. (2022).
\newblock Vision Transformers are robust learners.
\newblock In \textit{AAAI Conference on Artificial Intelligence}.

\bibitem{fort2023}
Fort, S., Ren, J., \& Lakshminarayanan, B. (2023).
\newblock Exploring the landscape of distributional robustness for vision transformers and convolutional neural networks.
\newblock In \textit{International Conference on Machine Learning}.

\end{thebibliography}

\end{document}
